\section{Introduction}\label{sec:intro}

Solving a linear system $Ax=b$ is the workhorse of scientific computing. Harrow, Hassidim and
Lloyd~\cite{hhl2009} observed that a quantum computer can prepare the state $\ket x\propto A^{-1}\ket
b$ for a Hermitian $A\in\C^{N\times N}$ in time polynomial in $\log N$, in the condition number
$\kappa$ and in $1/\varepsilon$, provided that $A$ is sparse and $\ket b$ is easy to prepare. The
algorithm, now called HHL, started the field of quantum linear algebra, and its fine print (state
preparation, sparse access, the condition number, and an output that is a quantum state rather than
a vector of numbers) has become part of how quantum speedups are judged~\cite{aaronson2015}.

HHL is also a demanding target for formal verification. Its specification is probabilistic: it
speaks about the outcome of a measurement and its probability, and the complete algorithm is a loop
whose number of rounds is random. Every step is approximate: eigenvalues are not exactly
representable in the phase-estimation register, $e^{iAt}$ has to be simulated, and the condition
number enters through rotation angles. The original paper controls these errors in a dense appendix
whose arguments are sketched rather than carried out, and earlier verifications, in
Coq~\cite{zhou2023coqq,wu2024local}, treat the exact case, in which every eigenvalue sits on the
phase-estimation grid, and take the controlled evolution, the state preparation and the rotation as
given unitaries.

\paragraph{What we verified.} This paper reports a verification of HHL in Isabelle/HOL that covers
the whole algorithm (Fig.~\ref{fig:stack}). Its headline is a single Hoare triple: in the exact
case, the repeat-until-success program
\[
\hoaret{I}{\ \mathit{hhl\_rus}\ }{\ketbra tt},\qquad \ket t=\ket0\otimes\frac{A^{-1}b}{\norm{A^{-1}b}}\otimes\ket1,
\]
terminates with probability one, and the system register then holds the normalised solution
(Theorem~\ref{thm:rus}). Below it lie the functional calculus of Hermitian matrices, phase
estimation and the error analysis of the original paper, including its Theorem~1 for arbitrary
spectra; above it, the programs are refined until every gate is a one- or two-qubit unitary. For
one concrete family, the shifted cycle Laplacian $L+\mu I$, the complete pipeline (state
preparation, Hamiltonian simulation by a Trotter product, two phase estimations, the rotation and
the flag) is a verified list of elementary gates of explicitly bounded, polynomial length
(Theorem~\ref{thm:lap}). On the way, the formalisation refuted seven statements of the HHL paper,
among them the constant of its Lemma~3 (Section~\ref{sec:errata}).

\begin{figure}[t]
\centering
\begin{tikzpicture}[font=\footnotesize,
  lay/.style={draw, rounded corners=2pt, text width=6.3cm, align=left, inner sep=3pt},
  when/.style={font=\scriptsize\itshape, text width=2.05cm, align=left, anchor=west},
  size/.style={font=\scriptsize, text width=1.75cm, align=right, anchor=east},
  arr/.style={-{Stealth[length=4pt]}, thick, gray!70}]
\node[lay, fill=notes!14] (L0) {\textbf{Foundation.} Paulsen's lecture course: spectral theorem, positivity,
  tensor products, measurement, channels, Choi--Kraus, error correction (\S\ref{sec:foundation})};
\node[lay, fill=hfc!14, above=0.3cm of L0] (L1) {\textbf{Mathematics.} Functional calculus of Hermitian matrices,
  Trotter and Suzuki product formulas, phase estimation (\S\ref{sec:math})};
\node[lay, fill=hhl!7, above=0.3cm of L1] (L2) {\textbf{Algorithm.} Exact HHL as a matrix identity; inexact
  phase estimation, filter functions, Theorem~1, run time (\S\ref{sec:algorithm})};
\node[lay, fill=hhl!12, above=0.3cm of L2] (L3) {\textbf{Program.} Quantum Hoare logic: one run, repeat until
  success, total correctness (\S\ref{sec:programs})};
\node[lay, fill=gates!12, above=0.3cm of L3] (L4) {\textbf{Circuit.} One- and two-qubit gates; $\ket b$, $e^{iAt}$
  and the rotation compiled; $L+\mu I$ with a gate count (\S\ref{sec:gates})};
\foreach \a/\b in {L0/L1, L1/L2, L2/L3, L3/L4} \draw[arr] (\a.north) -- (\b.south);
\node[when] at ($(L4.east)+(0.15,0)$) {Oct 21, 28--29};
\node[when] at ($(L3.east)+(0.15,0)$) {Oct 21--22};
\node[when] at ($(L2.east)+(0.15,0)$) {Oct 19--20};
\node[when] at ($(L1.east)+(0.15,0)$) {Oct 19--24};
\node[when] at ($(L0.east)+(0.15,0)$) {Apr 18--May 23,\\ Oct 18--20};
\node[size] at ($(L4.west)+(-0.12,0)$) {46 theories\\ 12.9k lines};
\node[size] at ($(L3.west)+(-0.12,0)$) {7 theories\\ 4.0k lines};
\node[size] at ($(L2.west)+(-0.12,0)$) {13 theories\\ 10.2k lines};
\node[size] at ($(L1.west)+(-0.12,0)$) {10 theories\\ 4.3k lines};
\node[size] at ($(L0.west)+(-0.12,0)$) {52 theories\\ 32.3k lines};
\end{tikzpicture}
\caption{The verified stack. Each level is proved from the one below; left, the size of each level in
the final development; right, the days on which it was built. The foundation was written in spring; everything above it was written after a pause of
almost five months, in eleven days of October 2026.}\label{fig:stack}
\end{figure}

\paragraph{How it was built.} The second subject of this paper is the process. The development was
written by an AI coding agent, Claude Code~\cite{claudecode}, under the direction of one
mathematician, the author, who wrote almost no proof text himself. It was not built in one sprint. A
formalisation of Vern Paulsen's lecture course on quantum computing and completely positive
maps~\cite{paulsen2011notes} was written in April and May 2026; then nothing happened for 147 days;
then, in eleven days of October, the course was completed and HHL was verified on top of it, from
the first theorem about phase estimation to the gate-level pipeline. Following the example of Ilin's
report on a Lean formalisation in plasma physics~\cite{ilin2026vml}, we report the timeline, the
statistics recoverable from the repository and the agent's logs, the role of the human, the ways in
which the agent failed, and what we learned. Our setting differs from his in three ways: the target
is a published algorithm whose paper turned out to contain errors, the work is split by a long
pause, and it builds on an earlier formalisation by the same team, which makes it a case study in
how an agent-built library ages and is reused.

\paragraph{Contributions.}
\begin{itemize}
\item A verification of HHL at four levels: mathematics, algorithm, program and circuit
(Sections~\ref{sec:hhl}--\ref{sec:results}). It includes total correctness of the repeat-until-success
loop, which needs a new proof rule for loops with a resetting body, the paper's error analysis with
explicit constants, and a gate-level pipeline with a closed-form gate count.
\item Errata to the HHL paper, each backed by a formal counterexample or a corrected theorem
(Section~\ref{sec:errata}).
\item A quantitative account of an agent-built formalisation that spans a long pause: 571 commits,
34 active days, 238 prompts in October, almost ten thousand tool calls, and the evolution of the
human's role from proving lemmas to judging statements, names and presentation
(Section~\ref{sec:process}).
\item Lessons on verifying quantum algorithms and on formalising with an agent, among them
the failure modes we observed and the checks that caught them (Section~\ref{sec:lessons}).
\end{itemize}

All results are checked by Isabelle2026 with \texttt{quick\_and\_dirty = false}, so no
\texttt{sorry} survives. An audit with \isa{Thm\_Deps.all\_oracles} over every theorem finds one
oracle, the Types-To-Sets rule \isa{cancel\_type\_definition}, and only in two side sessions, which
inherit it from imported AFP entries (Section~\ref{sec:trust}). The development is organised as
three entries for the Archive of Formal Proofs (AFP). The sources, the documents and the audits are
in the project repository;\footnote{\url{https://github.com/brydustin/Quantum_Computing}, commit
\texttt{1e37476}.} each listing below links to its source line (\isalogo).
